% ======================================================================
%  AT-MONO105 — Chapter 6: The D96 Spectrum
%  Canonical Monograph (v2.0) · The Actualization Theory:
%  A Reconstruction of Physics from Difference, Actualization and Spectrum
%
%  Status: COMPLETE — PASS-READY (MONO007 referee-ready architecture)
%  Inputs: Chapter01_Difference.tex, Chapter02_Eta.tex,
%          Chapter03_Actualization.tex, Chapter04_ClosureAndMinimalTheory.tex,
%          Chapter05_InevitableSpectrum.tex
%  Method: assembly only — canonical results only, no new physics, no speculation
%  Citations: QG155 (symmetry origin), QG156 (unified access law),
%             QG157 (effective access counts), QG158 (moment orders),
%             QG159 (D96 selection), QG295 (spectrum necessity)
%
%  Usage: % ======================================================================
%  AT-MONO105 — Chapter 6: The D96 Spectrum
%  Canonical Monograph (v2.0) · The Actualization Theory:
%  A Reconstruction of Physics from Difference, Actualization and Spectrum
%
%  Status: COMPLETE — PASS-READY (MONO007 referee-ready architecture)
%  Inputs: Chapter01_Difference.tex, Chapter02_Eta.tex,
%          Chapter03_Actualization.tex, Chapter04_ClosureAndMinimalTheory.tex,
%          Chapter05_InevitableSpectrum.tex
%  Method: assembly only — canonical results only, no new physics, no speculation
%  Citations: QG155 (symmetry origin), QG156 (unified access law),
%             QG157 (effective access counts), QG158 (moment orders),
%             QG159 (D96 selection), QG295 (spectrum necessity)
%
%  Usage: % ======================================================================
%  AT-MONO105 — Chapter 6: The D96 Spectrum
%  Canonical Monograph (v2.0) · The Actualization Theory:
%  A Reconstruction of Physics from Difference, Actualization and Spectrum
%
%  Status: COMPLETE — PASS-READY (MONO007 referee-ready architecture)
%  Inputs: Chapter01_Difference.tex, Chapter02_Eta.tex,
%          Chapter03_Actualization.tex, Chapter04_ClosureAndMinimalTheory.tex,
%          Chapter05_InevitableSpectrum.tex
%  Method: assembly only — canonical results only, no new physics, no speculation
%  Citations: QG155 (symmetry origin), QG156 (unified access law),
%             QG157 (effective access counts), QG158 (moment orders),
%             QG159 (D96 selection), QG295 (spectrum necessity)
%
%  Usage: % ======================================================================
%  AT-MONO105 — Chapter 6: The D96 Spectrum
%  Canonical Monograph (v2.0) · The Actualization Theory:
%  A Reconstruction of Physics from Difference, Actualization and Spectrum
%
%  Status: COMPLETE — PASS-READY (MONO007 referee-ready architecture)
%  Inputs: Chapter01_Difference.tex, Chapter02_Eta.tex,
%          Chapter03_Actualization.tex, Chapter04_ClosureAndMinimalTheory.tex,
%          Chapter05_InevitableSpectrum.tex
%  Method: assembly only — canonical results only, no new physics, no speculation
%  Citations: QG155 (symmetry origin), QG156 (unified access law),
%             QG157 (effective access counts), QG158 (moment orders),
%             QG159 (D96 selection), QG295 (spectrum necessity)
%
%  Usage: \input{Chapter06_D96Spectrum.tex} after the preamble
%         that defines the theorem environments (or include via the master file).
% ======================================================================

% ---- Theorem environments (define in the master preamble if not present) ----
% \newtheorem{definition}{Definition}[chapter]
% \newtheorem{lemma}[definition]{Lemma}
% \newtheorem{theorem}[definition]{Theorem}
% \newtheorem{corollary}[definition]{Corollary}
% \newtheorem{remark}[definition]{Remark}

\chapter{The D96 Spectrum}
\label{c06:chap:d96}

\section{Introduction}
\label{c06:sec:d96-introduction}

In Chapter~\ref{c05:chap:spectrum} we established that the spectrum of the theory is the
inevitable output of the actualization attractor: the unique Laplacian eigenspectrum of the
converged $N=96$ network, carrying no choice and presupposing no primitive status
(Theorem~\ref{c05:thm:spectrum-not-primitive}). This chapter
develops the concrete structure of that spectrum --- the \emph{D96 spectrum}.

We establish the structural content of the D96 spectrum in canonical form: the mode structure
(one zero mode, 95 positive modes); the spectral moments $\Sigma m$,
$\Sigma\sqrt{m}$, $\Sigma m^2$; the occupancy moment $\mathrm{occMom}$; the span; and the
multiplicity structure $[42\times2,\,5,\,6]$ \cite{c06:QG155,c06:QG156,c06:QG157,c06:QG158}. The D96 spectrum is
presented throughout as an \emph{emergent} structure --- the derived output of the
actualization attractor of Chapter~\ref{c05:chap:spectrum} --- never as a primitive: all
constants are derived from the spectrum, which is itself derived from the process
\cite{c06:QG159,c06:QG295}.

\section{The Spectrum Structure}
\label{c06:sec:spectrum-structure}

\begin{theorem}[The D96 spectrum is the derived output]
\label{c06:thm:d96-derived}
The \emph{D96 spectrum} is the Laplacian eigenspectrum of the converged $N=96$ attractor:
the set of $96$ eigenvalues of the graph Laplacian, consisting of one zero eigenvalue (the
background mode) and $95$ positive eigenvalues (the positive modes)
\cite{c06:QG155,c06:QG295}. It is the derived output of the actualization process, not a
primitive input: it admits no primitive interpretation and is not an underivable input.
\end{theorem}

\begin{proof}
By Theorem~\ref{c05:thm:spectrum-eigenspectrum} of Chapter~\ref{c05:chap:spectrum}, the spectrum is
the Laplacian eigenspectrum of the converged $N=96$ network, and by
Theorem~\ref{c05:thm:spectrum-not-primitive} of the same chapter it is not a primitive but the
inevitable output of the attractor. The spectrum-necessity audit \cite{c06:QG295} confirms
this classification explicitly --- the spectrum carries no choice and presupposes the
actualization process and its primitive Difference --- and no primitive interpretation of it
is consistent with the canonical architecture.
\end{proof}

\section{The Zero Mode}
\label{c06:sec:zero-mode}

\begin{theorem}[The zero mode is the background]
\label{c06:thm:zero-mode}
The \emph{zero mode} is the eigenvalue $\lambda_0 = 0$ of the graph Laplacian: the D96
spectrum has exactly one zero eigenvalue, corresponding to the background ---
the uniform configuration from which the count is measured.
\end{theorem}

\begin{proof}
The graph Laplacian of a connected network has exactly one zero eigenvalue, with the
constant eigenvector --- the uniform background, which the spectrum-necessity audit
\cite{c06:QG295} identifies as the background. By Definition~\ref{c01:def:difference} of
Chapter~\ref{c01:chap:difference}, Difference is the counting difference from a uniform
background; hence the zero mode is the reference against which Difference is measured, from
which the positive-mode structure represents the differences.
\end{proof}

\section{Positive Modes}
\label{c06:sec:positive-modes}

\begin{theorem}[There are 95 positive modes]
\label{c06:thm:95-modes}
The \emph{positive modes} are the $95$ positive eigenvalues
$\lambda_1,\dots,\lambda_{95}$ of the graph Laplacian: the non-background degrees of
freedom of the spectrum, which consists of exactly one zero mode and $95$ positive
eigenvalues.
\end{theorem}

\begin{proof}
The graph Laplacian of the $N=96$ network is a $96\times96$ matrix with one zero eigenvalue
(Theorem~\ref{c06:thm:zero-mode}) and the remaining $95$ eigenvalues positive; the
doublet-multiplicity structure confirms the count, $\Sigma m = 95$ being the total mode
count of the positive spectrum \cite{c06:QG157}, and the total dimension $96$ matches the
$N=96$ dimension of the attractor \cite{c06:QG155}.
\end{proof}

\section{Spectral Moments}
\label{c06:sec:spectral-moments}

\begin{theorem}[The moment ladder]
\label{c06:cor:moment-ladder}
The \emph{$p$-moment} of the multiplicity distribution is
\begin{equation}
\label{c06:eq:moment}
\Sigma m^p \;=\; \sum_i m_i^p,
\end{equation}
where the sum runs over the degenerate groups of the spectrum, with multiplicities $m_i$.
The spectral moments form a ladder of access counts: the first moment is the total mode
count, $\Sigma m = 95$; the half-moment is the neutral-sector access count,
$\Sigma\sqrt{m} = 64.08$; and the second moment is the doublet-occupancy access count,
$\Sigma m^2 = 229$ --- each an access count of a different sector of the theory
\cite{c06:QG157,c06:QG158}.
\end{theorem}

\begin{proof}
The effective-access-count program \cite{c06:QG157} establishes that the first moment of the
doublet-multiplicity distribution is the total mode count; by
Theorem~\ref{c06:thm:95-modes} this is $95$, so $\Sigma m = \sum_i m_i = 95$ --- the full
positive spectrum accessed uniformly. The same program establishes that the half-moment
$\Sigma\sqrt{m} = \sum_i \sqrt{m_i}$ is the neutral access count --- the statistical weight
each degenerate group contributes to the neutral channel --- and the second moment
$\Sigma m^2 = \sum_i m_i^2$ the doublet-occupancy access count. With the multiplicity
structure $[42\times2,5,6]$, $\Sigma\sqrt{m} = 64.08$ and
$\Sigma m^2 = 42\cdot 2^2 + 5^2 + 6^2 = 168+25+36 = 229$; the moment-orders program
\cite{c06:QG158} establishes that the moments form a Z2-power ladder, each corresponding to a
distinct sector access.
\end{proof}

\section{Spectral Constants}
\label{c06:sec:spectral-constants}

\begin{theorem}[The occupancy moment]
\label{c06:thm:occmom}
The \emph{octave-occupation moment} is
\begin{equation}
\label{c06:eq:occmom}
\mathrm{occMom} \;=\; \frac{\sum_{\text{octaves}} \mathrm{occ}_j^2}{\mathrm{occ}_0},
\end{equation}
where $\mathrm{occ}_j$ are the octave occupancies and $\mathrm{occ}_0$ the lightest-octave
occupancy. The octave-occupation moment of the D96 spectrum is
$\mathrm{occMom} = 1900.25$: the up-sector (dense-band) access count.
\end{theorem}

\begin{proof}
The effective-access-count program \cite{c06:QG157} establishes that the octave-occupation
moment $\mathrm{occMom} = \sum_j \mathrm{occ}_j^2/\mathrm{occ}_0$ is the dense-band access
count of the up sector; the octave occupancies of the D96 spectrum yield
$\mathrm{occMom} = 1900.25$.
\end{proof}

\begin{theorem}[The span]
\label{c06:thm:span}
The span of the D96 spectrum is
\begin{equation}
\label{c06:eq:span}
\mathrm{span} \;=\; \frac{\omega_{\max}}{\omega_{\min}} \;=\; 6.40,
\end{equation}
the ratio of the highest to the lowest positive frequency, with the modes
$\omega_k=\sqrt{\lambda_k}$ given by the Laplacian eigenvalues
$\lambda_k = 2\sum_{d=1}^{6}(1-\cos(2\pi dk/96))$ of the canonical attractor graph
$C_{96}(\pm1,\ldots,\pm6)$.
\end{theorem}

\begin{proof}
The span is the ratio of the extremal positive modes, $\omega_{\max}$ and $\omega_{\min}$;
the D96 spectrum yields the numerical value $\mathrm{span} = 6.40$, established by the
spectral-structure program \cite{c06:QG155}.
\end{proof}

\begin{corollary}[All constants are derived]
\label{c06:cor:constants-derived}
The spectral constants --- $\Sigma m=95$, $\Sigma\sqrt{m}=64.08$, $\Sigma m^2=229$,
$\mathrm{occMom}=1900.25$, $\mathrm{span}=6.40$ --- are all derived from the D96 spectrum,
which is itself derived from the attractor.
\end{corollary}

\begin{proof}
Each constant is a moment or structural feature of the spectrum
(Theorem~\ref{c06:cor:moment-ladder}, Theorem~\ref{c06:thm:occmom}, Theorem~\ref{c06:thm:span});
by Theorem~\ref{c06:thm:d96-derived}, the spectrum is the derived output of the attractor.
Hence every constant is twice derived --- a function of the spectrum, which is a function of
the process --- and no constant is a primitive input \cite{c06:QG157,c06:QG295}.
\end{proof}

\section{Multiplicity Structure}
\label{c06:sec:multiplicity}

\begin{theorem}[The multiplicity structure]
\label{c06:thm:multiplicity}
The \emph{multiplicity structure} of the D96 spectrum is the multiset of degenerate-group
sizes $m_i$: the pattern of equal eigenvalues grouped into the doublet structure,
$[42\times2,\,5,\,6]$ --- forty-two doublets, one group of size $5$, and one group of size
$6$. The doublet is the dominant multiplicity: $42$ of the $44$ groups are of size $2$,
reflecting the Z2 pairing of the spectrum generated by the D96 symmetry.
\end{theorem}

\begin{proof}
The doublet-multiplicity structure is established by the effective-access-count program
\cite{c06:QG157}: the spectrum is a multiset of degenerate groups with multiplicities $m_i$,
consisting of $42$ groups of size $2$, one of size $5$, and one of size $6$. The Z2 doublet
structure is generated by the D96 symmetry \cite{c06:QG155} --- the observable-sector
spectrum is fully Z2 paired --- and the total satisfies
$\Sigma m = 42\cdot2 + 5 + 6 = 95$, consistent with Theorem~\ref{c06:cor:moment-ladder}.
\end{proof}

\section{Consequences for Organization}
\label{c06:sec:organization}

The D96 spectrum is the substrate of the organization structure: its degeneracy groups,
moments, and span determine the organizational reads of the theory --- degeneracy (from the
multiplicity groups), compression and beat (from the span and octave structure), and
locking (from the spectral gap)
(Theorem~\ref{c06:thm:multiplicity}, Theorem~\ref{c06:thm:span}).

\begin{theorem}[The organization constants are derived]
\label{c06:thm:organization-derived}
The organization constants of the theory --- the occupancy moment, the span, and the
degeneracy structure --- are derived from the D96 spectrum, carrying no independent input.
The organization structure is therefore emergent: it is a read of the D96 spectrum and
carries no primitive status.
\end{theorem}

\begin{proof}
By Corollary~\ref{c06:cor:constants-derived}, the spectral constants are derived from the
spectrum, which is derived from the attractor; the organization constants are functions of
these spectral constants, so the organization structure inherits the derived status: every
organizational constant is an output of the spectrum, itself an output of the process; no
organizational feature is primitive \cite{c06:QG295}.
\end{proof}

Consistent with MONO007 \cite{c06:MONO007}, the D96 spectrum is presented throughout this
chapter as the derived output of the actualization attractor, never as a primitive: its
constants and multiplicities are all derived from the spectrum, which is itself derived.
This removes any possible reading of the D96 structure as an unexplained input.

\section{Summary}
\label{c06:sec:d96-summary}

This chapter has established the concrete structure of the inevitable spectrum. The D96
spectrum is the Laplacian eigenspectrum of the attractor (Theorem~\ref{c06:thm:d96-derived}):
one zero mode, the background that Difference is counted against
(Theorem~\ref{c06:thm:zero-mode}), and $95$ positive modes (Theorem~\ref{c06:thm:95-modes}).
Its moments $\Sigma m=95$, $\Sigma\sqrt{m}=64.08$, $\Sigma m^2=229$ form the access-count
ladder (Theorem~\ref{c06:cor:moment-ladder}); its constants $\mathrm{occMom}=1900.25$ and
$\mathrm{span}=6.40$ (Theorem~\ref{c06:thm:occmom}, Theorem~\ref{c06:thm:span}) and its
multiplicities $[42\times2,5,6]$ (Theorem~\ref{c06:thm:multiplicity}) are all derived from
the spectrum, which is itself derived from the actualization attractor
(Corollary~\ref{c06:cor:constants-derived}). The spectrum is therefore the substrate of the
organization structure (Theorem~\ref{c06:thm:organization-derived}); the following chapters
read the organization structure and then the physics sectors from this one spectrum.

\begin{thebibliography}{99}
\bibitem{c06:QG155} AT-QG Phase 155, \emph{Symmetry Origin} (Z2 doublet structure), Docs/Research.
\bibitem{c06:QG156} AT-QG Phase 156, \emph{Unified Spectral Access Law}, Docs/Research.
\bibitem{c06:QG157} AT-QG Phase 157, \emph{Effective Access Counts}, Docs/Research.
\bibitem{c06:QG158} AT-QG Phase 158, \emph{Moment Orders}, Docs/Research.
\bibitem{c06:QG159} AT-QG Phase 159, \emph{D96 Selection Origin}, Docs/Research.
\bibitem{c06:QG295} AT-QG Phase 295, \emph{Spectrum Necessity Audit}, Docs/Research.
\bibitem{c06:MONO007} AT-MONO007, \emph{Referee-Readiness Resolution}, Docs/Zenodo/Monograph.
\end{thebibliography}
 after the preamble
%         that defines the theorem environments (or include via the master file).
% ======================================================================

% ---- Theorem environments (define in the master preamble if not present) ----
% \newtheorem{definition}{Definition}[chapter]
% \newtheorem{lemma}[definition]{Lemma}
% \newtheorem{theorem}[definition]{Theorem}
% \newtheorem{corollary}[definition]{Corollary}
% \newtheorem{remark}[definition]{Remark}

\chapter{The D96 Spectrum}
\label{c06:chap:d96}

\section{Introduction}
\label{c06:sec:d96-introduction}

In Chapter~\ref{c05:chap:spectrum} we established that the spectrum of the theory is the
inevitable output of the actualization attractor: the unique Laplacian eigenspectrum of the
converged $N=96$ network, carrying no choice and presupposing no primitive status
(Theorem~\ref{c05:thm:spectrum-not-primitive}). This chapter
develops the concrete structure of that spectrum --- the \emph{D96 spectrum}.

We establish the structural content of the D96 spectrum in canonical form: the mode structure
(one zero mode, 95 positive modes); the spectral moments $\Sigma m$,
$\Sigma\sqrt{m}$, $\Sigma m^2$; the occupancy moment $\mathrm{occMom}$; the span; and the
multiplicity structure $[42\times2,\,5,\,6]$ \cite{c06:QG155,c06:QG156,c06:QG157,c06:QG158}. The D96 spectrum is
presented throughout as an \emph{emergent} structure --- the derived output of the
actualization attractor of Chapter~\ref{c05:chap:spectrum} --- never as a primitive: all
constants are derived from the spectrum, which is itself derived from the process
\cite{c06:QG159,c06:QG295}.

\section{The Spectrum Structure}
\label{c06:sec:spectrum-structure}

\begin{theorem}[The D96 spectrum is the derived output]
\label{c06:thm:d96-derived}
The \emph{D96 spectrum} is the Laplacian eigenspectrum of the converged $N=96$ attractor:
the set of $96$ eigenvalues of the graph Laplacian, consisting of one zero eigenvalue (the
background mode) and $95$ positive eigenvalues (the positive modes)
\cite{c06:QG155,c06:QG295}. It is the derived output of the actualization process, not a
primitive input: it admits no primitive interpretation and is not an underivable input.
\end{theorem}

\begin{proof}
By Theorem~\ref{c05:thm:spectrum-eigenspectrum} of Chapter~\ref{c05:chap:spectrum}, the spectrum is
the Laplacian eigenspectrum of the converged $N=96$ network, and by
Theorem~\ref{c05:thm:spectrum-not-primitive} of the same chapter it is not a primitive but the
inevitable output of the attractor. The spectrum-necessity audit \cite{c06:QG295} confirms
this classification explicitly --- the spectrum carries no choice and presupposes the
actualization process and its primitive Difference --- and no primitive interpretation of it
is consistent with the canonical architecture.
\end{proof}

\section{The Zero Mode}
\label{c06:sec:zero-mode}

\begin{theorem}[The zero mode is the background]
\label{c06:thm:zero-mode}
The \emph{zero mode} is the eigenvalue $\lambda_0 = 0$ of the graph Laplacian: the D96
spectrum has exactly one zero eigenvalue, corresponding to the background ---
the uniform configuration from which the count is measured.
\end{theorem}

\begin{proof}
The graph Laplacian of a connected network has exactly one zero eigenvalue, with the
constant eigenvector --- the uniform background, which the spectrum-necessity audit
\cite{c06:QG295} identifies as the background. By Definition~\ref{c01:def:difference} of
Chapter~\ref{c01:chap:difference}, Difference is the counting difference from a uniform
background; hence the zero mode is the reference against which Difference is measured, from
which the positive-mode structure represents the differences.
\end{proof}

\section{Positive Modes}
\label{c06:sec:positive-modes}

\begin{theorem}[There are 95 positive modes]
\label{c06:thm:95-modes}
The \emph{positive modes} are the $95$ positive eigenvalues
$\lambda_1,\dots,\lambda_{95}$ of the graph Laplacian: the non-background degrees of
freedom of the spectrum, which consists of exactly one zero mode and $95$ positive
eigenvalues.
\end{theorem}

\begin{proof}
The graph Laplacian of the $N=96$ network is a $96\times96$ matrix with one zero eigenvalue
(Theorem~\ref{c06:thm:zero-mode}) and the remaining $95$ eigenvalues positive; the
doublet-multiplicity structure confirms the count, $\Sigma m = 95$ being the total mode
count of the positive spectrum \cite{c06:QG157}, and the total dimension $96$ matches the
$N=96$ dimension of the attractor \cite{c06:QG155}.
\end{proof}

\section{Spectral Moments}
\label{c06:sec:spectral-moments}

\begin{theorem}[The moment ladder]
\label{c06:cor:moment-ladder}
The \emph{$p$-moment} of the multiplicity distribution is
\begin{equation}
\label{c06:eq:moment}
\Sigma m^p \;=\; \sum_i m_i^p,
\end{equation}
where the sum runs over the degenerate groups of the spectrum, with multiplicities $m_i$.
The spectral moments form a ladder of access counts: the first moment is the total mode
count, $\Sigma m = 95$; the half-moment is the neutral-sector access count,
$\Sigma\sqrt{m} = 64.08$; and the second moment is the doublet-occupancy access count,
$\Sigma m^2 = 229$ --- each an access count of a different sector of the theory
\cite{c06:QG157,c06:QG158}.
\end{theorem}

\begin{proof}
The effective-access-count program \cite{c06:QG157} establishes that the first moment of the
doublet-multiplicity distribution is the total mode count; by
Theorem~\ref{c06:thm:95-modes} this is $95$, so $\Sigma m = \sum_i m_i = 95$ --- the full
positive spectrum accessed uniformly. The same program establishes that the half-moment
$\Sigma\sqrt{m} = \sum_i \sqrt{m_i}$ is the neutral access count --- the statistical weight
each degenerate group contributes to the neutral channel --- and the second moment
$\Sigma m^2 = \sum_i m_i^2$ the doublet-occupancy access count. With the multiplicity
structure $[42\times2,5,6]$, $\Sigma\sqrt{m} = 64.08$ and
$\Sigma m^2 = 42\cdot 2^2 + 5^2 + 6^2 = 168+25+36 = 229$; the moment-orders program
\cite{c06:QG158} establishes that the moments form a Z2-power ladder, each corresponding to a
distinct sector access.
\end{proof}

\section{Spectral Constants}
\label{c06:sec:spectral-constants}

\begin{theorem}[The occupancy moment]
\label{c06:thm:occmom}
The \emph{octave-occupation moment} is
\begin{equation}
\label{c06:eq:occmom}
\mathrm{occMom} \;=\; \frac{\sum_{\text{octaves}} \mathrm{occ}_j^2}{\mathrm{occ}_0},
\end{equation}
where $\mathrm{occ}_j$ are the octave occupancies and $\mathrm{occ}_0$ the lightest-octave
occupancy. The octave-occupation moment of the D96 spectrum is
$\mathrm{occMom} = 1900.25$: the up-sector (dense-band) access count.
\end{theorem}

\begin{proof}
The effective-access-count program \cite{c06:QG157} establishes that the octave-occupation
moment $\mathrm{occMom} = \sum_j \mathrm{occ}_j^2/\mathrm{occ}_0$ is the dense-band access
count of the up sector; the octave occupancies of the D96 spectrum yield
$\mathrm{occMom} = 1900.25$.
\end{proof}

\begin{theorem}[The span]
\label{c06:thm:span}
The span of the D96 spectrum is
\begin{equation}
\label{c06:eq:span}
\mathrm{span} \;=\; \frac{\omega_{\max}}{\omega_{\min}} \;=\; 6.40,
\end{equation}
the ratio of the highest to the lowest positive frequency, with the modes
$\omega_k=\sqrt{\lambda_k}$ given by the Laplacian eigenvalues
$\lambda_k = 2\sum_{d=1}^{6}(1-\cos(2\pi dk/96))$ of the canonical attractor graph
$C_{96}(\pm1,\ldots,\pm6)$.
\end{theorem}

\begin{proof}
The span is the ratio of the extremal positive modes, $\omega_{\max}$ and $\omega_{\min}$;
the D96 spectrum yields the numerical value $\mathrm{span} = 6.40$, established by the
spectral-structure program \cite{c06:QG155}.
\end{proof}

\begin{corollary}[All constants are derived]
\label{c06:cor:constants-derived}
The spectral constants --- $\Sigma m=95$, $\Sigma\sqrt{m}=64.08$, $\Sigma m^2=229$,
$\mathrm{occMom}=1900.25$, $\mathrm{span}=6.40$ --- are all derived from the D96 spectrum,
which is itself derived from the attractor.
\end{corollary}

\begin{proof}
Each constant is a moment or structural feature of the spectrum
(Theorem~\ref{c06:cor:moment-ladder}, Theorem~\ref{c06:thm:occmom}, Theorem~\ref{c06:thm:span});
by Theorem~\ref{c06:thm:d96-derived}, the spectrum is the derived output of the attractor.
Hence every constant is twice derived --- a function of the spectrum, which is a function of
the process --- and no constant is a primitive input \cite{c06:QG157,c06:QG295}.
\end{proof}

\section{Multiplicity Structure}
\label{c06:sec:multiplicity}

\begin{theorem}[The multiplicity structure]
\label{c06:thm:multiplicity}
The \emph{multiplicity structure} of the D96 spectrum is the multiset of degenerate-group
sizes $m_i$: the pattern of equal eigenvalues grouped into the doublet structure,
$[42\times2,\,5,\,6]$ --- forty-two doublets, one group of size $5$, and one group of size
$6$. The doublet is the dominant multiplicity: $42$ of the $44$ groups are of size $2$,
reflecting the Z2 pairing of the spectrum generated by the D96 symmetry.
\end{theorem}

\begin{proof}
The doublet-multiplicity structure is established by the effective-access-count program
\cite{c06:QG157}: the spectrum is a multiset of degenerate groups with multiplicities $m_i$,
consisting of $42$ groups of size $2$, one of size $5$, and one of size $6$. The Z2 doublet
structure is generated by the D96 symmetry \cite{c06:QG155} --- the observable-sector
spectrum is fully Z2 paired --- and the total satisfies
$\Sigma m = 42\cdot2 + 5 + 6 = 95$, consistent with Theorem~\ref{c06:cor:moment-ladder}.
\end{proof}

\section{Consequences for Organization}
\label{c06:sec:organization}

The D96 spectrum is the substrate of the organization structure: its degeneracy groups,
moments, and span determine the organizational reads of the theory --- degeneracy (from the
multiplicity groups), compression and beat (from the span and octave structure), and
locking (from the spectral gap)
(Theorem~\ref{c06:thm:multiplicity}, Theorem~\ref{c06:thm:span}).

\begin{theorem}[The organization constants are derived]
\label{c06:thm:organization-derived}
The organization constants of the theory --- the occupancy moment, the span, and the
degeneracy structure --- are derived from the D96 spectrum, carrying no independent input.
The organization structure is therefore emergent: it is a read of the D96 spectrum and
carries no primitive status.
\end{theorem}

\begin{proof}
By Corollary~\ref{c06:cor:constants-derived}, the spectral constants are derived from the
spectrum, which is derived from the attractor; the organization constants are functions of
these spectral constants, so the organization structure inherits the derived status: every
organizational constant is an output of the spectrum, itself an output of the process; no
organizational feature is primitive \cite{c06:QG295}.
\end{proof}

Consistent with MONO007 \cite{c06:MONO007}, the D96 spectrum is presented throughout this
chapter as the derived output of the actualization attractor, never as a primitive: its
constants and multiplicities are all derived from the spectrum, which is itself derived.
This removes any possible reading of the D96 structure as an unexplained input.

\section{Summary}
\label{c06:sec:d96-summary}

This chapter has established the concrete structure of the inevitable spectrum. The D96
spectrum is the Laplacian eigenspectrum of the attractor (Theorem~\ref{c06:thm:d96-derived}):
one zero mode, the background that Difference is counted against
(Theorem~\ref{c06:thm:zero-mode}), and $95$ positive modes (Theorem~\ref{c06:thm:95-modes}).
Its moments $\Sigma m=95$, $\Sigma\sqrt{m}=64.08$, $\Sigma m^2=229$ form the access-count
ladder (Theorem~\ref{c06:cor:moment-ladder}); its constants $\mathrm{occMom}=1900.25$ and
$\mathrm{span}=6.40$ (Theorem~\ref{c06:thm:occmom}, Theorem~\ref{c06:thm:span}) and its
multiplicities $[42\times2,5,6]$ (Theorem~\ref{c06:thm:multiplicity}) are all derived from
the spectrum, which is itself derived from the actualization attractor
(Corollary~\ref{c06:cor:constants-derived}). The spectrum is therefore the substrate of the
organization structure (Theorem~\ref{c06:thm:organization-derived}); the following chapters
read the organization structure and then the physics sectors from this one spectrum.

\begin{thebibliography}{99}
\bibitem{c06:QG155} AT-QG Phase 155, \emph{Symmetry Origin} (Z2 doublet structure), Docs/Research.
\bibitem{c06:QG156} AT-QG Phase 156, \emph{Unified Spectral Access Law}, Docs/Research.
\bibitem{c06:QG157} AT-QG Phase 157, \emph{Effective Access Counts}, Docs/Research.
\bibitem{c06:QG158} AT-QG Phase 158, \emph{Moment Orders}, Docs/Research.
\bibitem{c06:QG159} AT-QG Phase 159, \emph{D96 Selection Origin}, Docs/Research.
\bibitem{c06:QG295} AT-QG Phase 295, \emph{Spectrum Necessity Audit}, Docs/Research.
\bibitem{c06:MONO007} AT-MONO007, \emph{Referee-Readiness Resolution}, Docs/Zenodo/Monograph.
\end{thebibliography}
 after the preamble
%         that defines the theorem environments (or include via the master file).
% ======================================================================

% ---- Theorem environments (define in the master preamble if not present) ----
% \newtheorem{definition}{Definition}[chapter]
% \newtheorem{lemma}[definition]{Lemma}
% \newtheorem{theorem}[definition]{Theorem}
% \newtheorem{corollary}[definition]{Corollary}
% \newtheorem{remark}[definition]{Remark}

\chapter{The D96 Spectrum}
\label{c06:chap:d96}

\section{Introduction}
\label{c06:sec:d96-introduction}

In Chapter~\ref{c05:chap:spectrum} we established that the spectrum of the theory is the
inevitable output of the actualization attractor: the unique Laplacian eigenspectrum of the
converged $N=96$ network, carrying no choice and presupposing no primitive status
(Theorem~\ref{c05:thm:spectrum-not-primitive}). This chapter
develops the concrete structure of that spectrum --- the \emph{D96 spectrum}.

We establish the structural content of the D96 spectrum in canonical form: the mode structure
(one zero mode, 95 positive modes); the spectral moments $\Sigma m$,
$\Sigma\sqrt{m}$, $\Sigma m^2$; the occupancy moment $\mathrm{occMom}$; the span; and the
multiplicity structure $[42\times2,\,5,\,6]$ \cite{c06:QG155,c06:QG156,c06:QG157,c06:QG158}. The D96 spectrum is
presented throughout as an \emph{emergent} structure --- the derived output of the
actualization attractor of Chapter~\ref{c05:chap:spectrum} --- never as a primitive: all
constants are derived from the spectrum, which is itself derived from the process
\cite{c06:QG159,c06:QG295}.

\section{The Spectrum Structure}
\label{c06:sec:spectrum-structure}

\begin{theorem}[The D96 spectrum is the derived output]
\label{c06:thm:d96-derived}
The \emph{D96 spectrum} is the Laplacian eigenspectrum of the converged $N=96$ attractor:
the set of $96$ eigenvalues of the graph Laplacian, consisting of one zero eigenvalue (the
background mode) and $95$ positive eigenvalues (the positive modes)
\cite{c06:QG155,c06:QG295}. It is the derived output of the actualization process, not a
primitive input: it admits no primitive interpretation and is not an underivable input.
\end{theorem}

\begin{proof}
By Theorem~\ref{c05:thm:spectrum-eigenspectrum} of Chapter~\ref{c05:chap:spectrum}, the spectrum is
the Laplacian eigenspectrum of the converged $N=96$ network, and by
Theorem~\ref{c05:thm:spectrum-not-primitive} of the same chapter it is not a primitive but the
inevitable output of the attractor. The spectrum-necessity audit \cite{c06:QG295} confirms
this classification explicitly --- the spectrum carries no choice and presupposes the
actualization process and its primitive Difference --- and no primitive interpretation of it
is consistent with the canonical architecture.
\end{proof}

\section{The Zero Mode}
\label{c06:sec:zero-mode}

\begin{theorem}[The zero mode is the background]
\label{c06:thm:zero-mode}
The \emph{zero mode} is the eigenvalue $\lambda_0 = 0$ of the graph Laplacian: the D96
spectrum has exactly one zero eigenvalue, corresponding to the background ---
the uniform configuration from which the count is measured.
\end{theorem}

\begin{proof}
The graph Laplacian of a connected network has exactly one zero eigenvalue, with the
constant eigenvector --- the uniform background, which the spectrum-necessity audit
\cite{c06:QG295} identifies as the background. By Definition~\ref{c01:def:difference} of
Chapter~\ref{c01:chap:difference}, Difference is the counting difference from a uniform
background; hence the zero mode is the reference against which Difference is measured, from
which the positive-mode structure represents the differences.
\end{proof}

\section{Positive Modes}
\label{c06:sec:positive-modes}

\begin{theorem}[There are 95 positive modes]
\label{c06:thm:95-modes}
The \emph{positive modes} are the $95$ positive eigenvalues
$\lambda_1,\dots,\lambda_{95}$ of the graph Laplacian: the non-background degrees of
freedom of the spectrum, which consists of exactly one zero mode and $95$ positive
eigenvalues.
\end{theorem}

\begin{proof}
The graph Laplacian of the $N=96$ network is a $96\times96$ matrix with one zero eigenvalue
(Theorem~\ref{c06:thm:zero-mode}) and the remaining $95$ eigenvalues positive; the
doublet-multiplicity structure confirms the count, $\Sigma m = 95$ being the total mode
count of the positive spectrum \cite{c06:QG157}, and the total dimension $96$ matches the
$N=96$ dimension of the attractor \cite{c06:QG155}.
\end{proof}

\section{Spectral Moments}
\label{c06:sec:spectral-moments}

\begin{theorem}[The moment ladder]
\label{c06:cor:moment-ladder}
The \emph{$p$-moment} of the multiplicity distribution is
\begin{equation}
\label{c06:eq:moment}
\Sigma m^p \;=\; \sum_i m_i^p,
\end{equation}
where the sum runs over the degenerate groups of the spectrum, with multiplicities $m_i$.
The spectral moments form a ladder of access counts: the first moment is the total mode
count, $\Sigma m = 95$; the half-moment is the neutral-sector access count,
$\Sigma\sqrt{m} = 64.08$; and the second moment is the doublet-occupancy access count,
$\Sigma m^2 = 229$ --- each an access count of a different sector of the theory
\cite{c06:QG157,c06:QG158}.
\end{theorem}

\begin{proof}
The effective-access-count program \cite{c06:QG157} establishes that the first moment of the
doublet-multiplicity distribution is the total mode count; by
Theorem~\ref{c06:thm:95-modes} this is $95$, so $\Sigma m = \sum_i m_i = 95$ --- the full
positive spectrum accessed uniformly. The same program establishes that the half-moment
$\Sigma\sqrt{m} = \sum_i \sqrt{m_i}$ is the neutral access count --- the statistical weight
each degenerate group contributes to the neutral channel --- and the second moment
$\Sigma m^2 = \sum_i m_i^2$ the doublet-occupancy access count. With the multiplicity
structure $[42\times2,5,6]$, $\Sigma\sqrt{m} = 64.08$ and
$\Sigma m^2 = 42\cdot 2^2 + 5^2 + 6^2 = 168+25+36 = 229$; the moment-orders program
\cite{c06:QG158} establishes that the moments form a Z2-power ladder, each corresponding to a
distinct sector access.
\end{proof}

\section{Spectral Constants}
\label{c06:sec:spectral-constants}

\begin{theorem}[The occupancy moment]
\label{c06:thm:occmom}
The \emph{octave-occupation moment} is
\begin{equation}
\label{c06:eq:occmom}
\mathrm{occMom} \;=\; \frac{\sum_{\text{octaves}} \mathrm{occ}_j^2}{\mathrm{occ}_0},
\end{equation}
where $\mathrm{occ}_j$ are the octave occupancies and $\mathrm{occ}_0$ the lightest-octave
occupancy. The octave-occupation moment of the D96 spectrum is
$\mathrm{occMom} = 1900.25$: the up-sector (dense-band) access count.
\end{theorem}

\begin{proof}
The effective-access-count program \cite{c06:QG157} establishes that the octave-occupation
moment $\mathrm{occMom} = \sum_j \mathrm{occ}_j^2/\mathrm{occ}_0$ is the dense-band access
count of the up sector; the octave occupancies of the D96 spectrum yield
$\mathrm{occMom} = 1900.25$.
\end{proof}

\begin{theorem}[The span]
\label{c06:thm:span}
The span of the D96 spectrum is
\begin{equation}
\label{c06:eq:span}
\mathrm{span} \;=\; \frac{\omega_{\max}}{\omega_{\min}} \;=\; 6.40,
\end{equation}
the ratio of the highest to the lowest positive frequency, with the modes
$\omega_k=\sqrt{\lambda_k}$ given by the Laplacian eigenvalues
$\lambda_k = 2\sum_{d=1}^{6}(1-\cos(2\pi dk/96))$ of the canonical attractor graph
$C_{96}(\pm1,\ldots,\pm6)$.
\end{theorem}

\begin{proof}
The span is the ratio of the extremal positive modes, $\omega_{\max}$ and $\omega_{\min}$;
the D96 spectrum yields the numerical value $\mathrm{span} = 6.40$, established by the
spectral-structure program \cite{c06:QG155}.
\end{proof}

\begin{corollary}[All constants are derived]
\label{c06:cor:constants-derived}
The spectral constants --- $\Sigma m=95$, $\Sigma\sqrt{m}=64.08$, $\Sigma m^2=229$,
$\mathrm{occMom}=1900.25$, $\mathrm{span}=6.40$ --- are all derived from the D96 spectrum,
which is itself derived from the attractor.
\end{corollary}

\begin{proof}
Each constant is a moment or structural feature of the spectrum
(Theorem~\ref{c06:cor:moment-ladder}, Theorem~\ref{c06:thm:occmom}, Theorem~\ref{c06:thm:span});
by Theorem~\ref{c06:thm:d96-derived}, the spectrum is the derived output of the attractor.
Hence every constant is twice derived --- a function of the spectrum, which is a function of
the process --- and no constant is a primitive input \cite{c06:QG157,c06:QG295}.
\end{proof}

\section{Multiplicity Structure}
\label{c06:sec:multiplicity}

\begin{theorem}[The multiplicity structure]
\label{c06:thm:multiplicity}
The \emph{multiplicity structure} of the D96 spectrum is the multiset of degenerate-group
sizes $m_i$: the pattern of equal eigenvalues grouped into the doublet structure,
$[42\times2,\,5,\,6]$ --- forty-two doublets, one group of size $5$, and one group of size
$6$. The doublet is the dominant multiplicity: $42$ of the $44$ groups are of size $2$,
reflecting the Z2 pairing of the spectrum generated by the D96 symmetry.
\end{theorem}

\begin{proof}
The doublet-multiplicity structure is established by the effective-access-count program
\cite{c06:QG157}: the spectrum is a multiset of degenerate groups with multiplicities $m_i$,
consisting of $42$ groups of size $2$, one of size $5$, and one of size $6$. The Z2 doublet
structure is generated by the D96 symmetry \cite{c06:QG155} --- the observable-sector
spectrum is fully Z2 paired --- and the total satisfies
$\Sigma m = 42\cdot2 + 5 + 6 = 95$, consistent with Theorem~\ref{c06:cor:moment-ladder}.
\end{proof}

\section{Consequences for Organization}
\label{c06:sec:organization}

The D96 spectrum is the substrate of the organization structure: its degeneracy groups,
moments, and span determine the organizational reads of the theory --- degeneracy (from the
multiplicity groups), compression and beat (from the span and octave structure), and
locking (from the spectral gap)
(Theorem~\ref{c06:thm:multiplicity}, Theorem~\ref{c06:thm:span}).

\begin{theorem}[The organization constants are derived]
\label{c06:thm:organization-derived}
The organization constants of the theory --- the occupancy moment, the span, and the
degeneracy structure --- are derived from the D96 spectrum, carrying no independent input.
The organization structure is therefore emergent: it is a read of the D96 spectrum and
carries no primitive status.
\end{theorem}

\begin{proof}
By Corollary~\ref{c06:cor:constants-derived}, the spectral constants are derived from the
spectrum, which is derived from the attractor; the organization constants are functions of
these spectral constants, so the organization structure inherits the derived status: every
organizational constant is an output of the spectrum, itself an output of the process; no
organizational feature is primitive \cite{c06:QG295}.
\end{proof}

Consistent with MONO007 \cite{c06:MONO007}, the D96 spectrum is presented throughout this
chapter as the derived output of the actualization attractor, never as a primitive: its
constants and multiplicities are all derived from the spectrum, which is itself derived.
This removes any possible reading of the D96 structure as an unexplained input.

\section{Summary}
\label{c06:sec:d96-summary}

This chapter has established the concrete structure of the inevitable spectrum. The D96
spectrum is the Laplacian eigenspectrum of the attractor (Theorem~\ref{c06:thm:d96-derived}):
one zero mode, the background that Difference is counted against
(Theorem~\ref{c06:thm:zero-mode}), and $95$ positive modes (Theorem~\ref{c06:thm:95-modes}).
Its moments $\Sigma m=95$, $\Sigma\sqrt{m}=64.08$, $\Sigma m^2=229$ form the access-count
ladder (Theorem~\ref{c06:cor:moment-ladder}); its constants $\mathrm{occMom}=1900.25$ and
$\mathrm{span}=6.40$ (Theorem~\ref{c06:thm:occmom}, Theorem~\ref{c06:thm:span}) and its
multiplicities $[42\times2,5,6]$ (Theorem~\ref{c06:thm:multiplicity}) are all derived from
the spectrum, which is itself derived from the actualization attractor
(Corollary~\ref{c06:cor:constants-derived}). The spectrum is therefore the substrate of the
organization structure (Theorem~\ref{c06:thm:organization-derived}); the following chapters
read the organization structure and then the physics sectors from this one spectrum.

\begin{thebibliography}{99}
\bibitem{c06:QG155} AT-QG Phase 155, \emph{Symmetry Origin} (Z2 doublet structure), Docs/Research.
\bibitem{c06:QG156} AT-QG Phase 156, \emph{Unified Spectral Access Law}, Docs/Research.
\bibitem{c06:QG157} AT-QG Phase 157, \emph{Effective Access Counts}, Docs/Research.
\bibitem{c06:QG158} AT-QG Phase 158, \emph{Moment Orders}, Docs/Research.
\bibitem{c06:QG159} AT-QG Phase 159, \emph{D96 Selection Origin}, Docs/Research.
\bibitem{c06:QG295} AT-QG Phase 295, \emph{Spectrum Necessity Audit}, Docs/Research.
\bibitem{c06:MONO007} AT-MONO007, \emph{Referee-Readiness Resolution}, Docs/Zenodo/Monograph.
\end{thebibliography}
 after the preamble
%         that defines the theorem environments (or include via the master file).
% ======================================================================

% ---- Theorem environments (define in the master preamble if not present) ----
% \newtheorem{definition}{Definition}[chapter]
% \newtheorem{lemma}[definition]{Lemma}
% \newtheorem{theorem}[definition]{Theorem}
% \newtheorem{corollary}[definition]{Corollary}
% \newtheorem{remark}[definition]{Remark}

\chapter{The D96 Spectrum}
\label{c06:chap:d96}

\section{Introduction}
\label{c06:sec:d96-introduction}

In Chapter~\ref{c05:chap:spectrum} we established that the spectrum of the theory is the
inevitable output of the actualization attractor: the unique Laplacian eigenspectrum of the
converged $N=96$ network, carrying no choice and presupposing no primitive status
(Theorem~\ref{c05:thm:spectrum-not-primitive}). This chapter
develops the concrete structure of that spectrum --- the \emph{D96 spectrum}.

We establish the structural content of the D96 spectrum in canonical form: the mode structure
(one zero mode, 95 positive modes); the spectral moments $\Sigma m$,
$\Sigma\sqrt{m}$, $\Sigma m^2$; the occupancy moment $\mathrm{occMom}$; the span; and the
multiplicity structure $[42\times2,\,5,\,6]$ \cite{c06:QG155,c06:QG156,c06:QG157,c06:QG158}. The D96 spectrum is
presented throughout as an \emph{emergent} structure --- the derived output of the
actualization attractor of Chapter~\ref{c05:chap:spectrum} --- never as a primitive: all
constants are derived from the spectrum, which is itself derived from the process
\cite{c06:QG159,c06:QG295}.

\section{The Spectrum Structure}
\label{c06:sec:spectrum-structure}

\begin{theorem}[The D96 spectrum is the derived output]
\label{c06:thm:d96-derived}
The \emph{D96 spectrum} is the Laplacian eigenspectrum of the converged $N=96$ attractor:
the set of $96$ eigenvalues of the graph Laplacian, consisting of one zero eigenvalue (the
background mode) and $95$ positive eigenvalues (the positive modes)
\cite{c06:QG155,c06:QG295}. It is the derived output of the actualization process, not a
primitive input: it admits no primitive interpretation and is not an underivable input.
\end{theorem}

\begin{proof}
By Theorem~\ref{c05:thm:spectrum-eigenspectrum} of Chapter~\ref{c05:chap:spectrum}, the spectrum is
the Laplacian eigenspectrum of the converged $N=96$ network, and by
Theorem~\ref{c05:thm:spectrum-not-primitive} of the same chapter it is not a primitive but the
inevitable output of the attractor. The spectrum-necessity audit \cite{c06:QG295} confirms
this classification explicitly --- the spectrum carries no choice and presupposes the
actualization process and its primitive Difference --- and no primitive interpretation of it
is consistent with the canonical architecture.
\end{proof}

\section{The Zero Mode}
\label{c06:sec:zero-mode}

\begin{theorem}[The zero mode is the background]
\label{c06:thm:zero-mode}
The \emph{zero mode} is the eigenvalue $\lambda_0 = 0$ of the graph Laplacian: the D96
spectrum has exactly one zero eigenvalue, corresponding to the background ---
the uniform configuration from which the count is measured.
\end{theorem}

\begin{proof}
The graph Laplacian of a connected network has exactly one zero eigenvalue, with the
constant eigenvector --- the uniform background, which the spectrum-necessity audit
\cite{c06:QG295} identifies as the background. By Definition~\ref{c01:def:difference} of
Chapter~\ref{c01:chap:difference}, Difference is the counting difference from a uniform
background; hence the zero mode is the reference against which Difference is measured, from
which the positive-mode structure represents the differences.
\end{proof}

\section{Positive Modes}
\label{c06:sec:positive-modes}

\begin{theorem}[There are 95 positive modes]
\label{c06:thm:95-modes}
The \emph{positive modes} are the $95$ positive eigenvalues
$\lambda_1,\dots,\lambda_{95}$ of the graph Laplacian: the non-background degrees of
freedom of the spectrum, which consists of exactly one zero mode and $95$ positive
eigenvalues.
\end{theorem}

\begin{proof}
The graph Laplacian of the $N=96$ network is a $96\times96$ matrix with one zero eigenvalue
(Theorem~\ref{c06:thm:zero-mode}) and the remaining $95$ eigenvalues positive; the
doublet-multiplicity structure confirms the count, $\Sigma m = 95$ being the total mode
count of the positive spectrum \cite{c06:QG157}, and the total dimension $96$ matches the
$N=96$ dimension of the attractor \cite{c06:QG155}.
\end{proof}

\section{Spectral Moments}
\label{c06:sec:spectral-moments}

\begin{theorem}[The moment ladder]
\label{c06:cor:moment-ladder}
The \emph{$p$-moment} of the multiplicity distribution is
\begin{equation}
\label{c06:eq:moment}
\Sigma m^p \;=\; \sum_i m_i^p,
\end{equation}
where the sum runs over the degenerate groups of the spectrum, with multiplicities $m_i$.
The spectral moments form a ladder of access counts: the first moment is the total mode
count, $\Sigma m = 95$; the half-moment is the neutral-sector access count,
$\Sigma\sqrt{m} = 64.08$; and the second moment is the doublet-occupancy access count,
$\Sigma m^2 = 229$ --- each an access count of a different sector of the theory
\cite{c06:QG157,c06:QG158}.
\end{theorem}

\begin{proof}
The effective-access-count program \cite{c06:QG157} establishes that the first moment of the
doublet-multiplicity distribution is the total mode count; by
Theorem~\ref{c06:thm:95-modes} this is $95$, so $\Sigma m = \sum_i m_i = 95$ --- the full
positive spectrum accessed uniformly. The same program establishes that the half-moment
$\Sigma\sqrt{m} = \sum_i \sqrt{m_i}$ is the neutral access count --- the statistical weight
each degenerate group contributes to the neutral channel --- and the second moment
$\Sigma m^2 = \sum_i m_i^2$ the doublet-occupancy access count. With the multiplicity
structure $[42\times2,5,6]$, $\Sigma\sqrt{m} = 64.08$ and
$\Sigma m^2 = 42\cdot 2^2 + 5^2 + 6^2 = 168+25+36 = 229$; the moment-orders program
\cite{c06:QG158} establishes that the moments form a Z2-power ladder, each corresponding to a
distinct sector access.
\end{proof}

\section{Spectral Constants}
\label{c06:sec:spectral-constants}

\begin{theorem}[The occupancy moment]
\label{c06:thm:occmom}
The \emph{octave-occupation moment} is
\begin{equation}
\label{c06:eq:occmom}
\mathrm{occMom} \;=\; \frac{\sum_{\text{octaves}} \mathrm{occ}_j^2}{\mathrm{occ}_0},
\end{equation}
where $\mathrm{occ}_j$ are the octave occupancies and $\mathrm{occ}_0$ the lightest-octave
occupancy. The octave-occupation moment of the D96 spectrum is
$\mathrm{occMom} = 1900.25$: the up-sector (dense-band) access count.
\end{theorem}

\begin{proof}
The effective-access-count program \cite{c06:QG157} establishes that the octave-occupation
moment $\mathrm{occMom} = \sum_j \mathrm{occ}_j^2/\mathrm{occ}_0$ is the dense-band access
count of the up sector; the octave occupancies of the D96 spectrum yield
$\mathrm{occMom} = 1900.25$.
\end{proof}

\begin{theorem}[The span]
\label{c06:thm:span}
The span of the D96 spectrum is
\begin{equation}
\label{c06:eq:span}
\mathrm{span} \;=\; \frac{\omega_{\max}}{\omega_{\min}} \;=\; 6.40,
\end{equation}
the ratio of the highest to the lowest positive frequency, with the modes
$\omega_k=\sqrt{\lambda_k}$ given by the Laplacian eigenvalues
$\lambda_k = 2\sum_{d=1}^{6}(1-\cos(2\pi dk/96))$ of the canonical attractor graph
$C_{96}(\pm1,\ldots,\pm6)$.
\end{theorem}

\begin{proof}
The span is the ratio of the extremal positive modes, $\omega_{\max}$ and $\omega_{\min}$;
the D96 spectrum yields the numerical value $\mathrm{span} = 6.40$, established by the
spectral-structure program \cite{c06:QG155}.
\end{proof}

\begin{corollary}[All constants are derived]
\label{c06:cor:constants-derived}
The spectral constants --- $\Sigma m=95$, $\Sigma\sqrt{m}=64.08$, $\Sigma m^2=229$,
$\mathrm{occMom}=1900.25$, $\mathrm{span}=6.40$ --- are all derived from the D96 spectrum,
which is itself derived from the attractor.
\end{corollary}

\begin{proof}
Each constant is a moment or structural feature of the spectrum
(Theorem~\ref{c06:cor:moment-ladder}, Theorem~\ref{c06:thm:occmom}, Theorem~\ref{c06:thm:span});
by Theorem~\ref{c06:thm:d96-derived}, the spectrum is the derived output of the attractor.
Hence every constant is twice derived --- a function of the spectrum, which is a function of
the process --- and no constant is a primitive input \cite{c06:QG157,c06:QG295}.
\end{proof}

\section{Multiplicity Structure}
\label{c06:sec:multiplicity}

\begin{theorem}[The multiplicity structure]
\label{c06:thm:multiplicity}
The \emph{multiplicity structure} of the D96 spectrum is the multiset of degenerate-group
sizes $m_i$: the pattern of equal eigenvalues grouped into the doublet structure,
$[42\times2,\,5,\,6]$ --- forty-two doublets, one group of size $5$, and one group of size
$6$. The doublet is the dominant multiplicity: $42$ of the $44$ groups are of size $2$,
reflecting the Z2 pairing of the spectrum generated by the D96 symmetry.
\end{theorem}

\begin{proof}
The doublet-multiplicity structure is established by the effective-access-count program
\cite{c06:QG157}: the spectrum is a multiset of degenerate groups with multiplicities $m_i$,
consisting of $42$ groups of size $2$, one of size $5$, and one of size $6$. The Z2 doublet
structure is generated by the D96 symmetry \cite{c06:QG155} --- the observable-sector
spectrum is fully Z2 paired --- and the total satisfies
$\Sigma m = 42\cdot2 + 5 + 6 = 95$, consistent with Theorem~\ref{c06:cor:moment-ladder}.
\end{proof}

\section{Consequences for Organization}
\label{c06:sec:organization}

The D96 spectrum is the substrate of the organization structure: its degeneracy groups,
moments, and span determine the organizational reads of the theory --- degeneracy (from the
multiplicity groups), compression and beat (from the span and octave structure), and
locking (from the spectral gap)
(Theorem~\ref{c06:thm:multiplicity}, Theorem~\ref{c06:thm:span}).

\begin{theorem}[The organization constants are derived]
\label{c06:thm:organization-derived}
The organization constants of the theory --- the occupancy moment, the span, and the
degeneracy structure --- are derived from the D96 spectrum, carrying no independent input.
The organization structure is therefore emergent: it is a read of the D96 spectrum and
carries no primitive status.
\end{theorem}

\begin{proof}
By Corollary~\ref{c06:cor:constants-derived}, the spectral constants are derived from the
spectrum, which is derived from the attractor; the organization constants are functions of
these spectral constants, so the organization structure inherits the derived status: every
organizational constant is an output of the spectrum, itself an output of the process; no
organizational feature is primitive \cite{c06:QG295}.
\end{proof}

Consistent with MONO007 \cite{c06:MONO007}, the D96 spectrum is presented throughout this
chapter as the derived output of the actualization attractor, never as a primitive: its
constants and multiplicities are all derived from the spectrum, which is itself derived.
This removes any possible reading of the D96 structure as an unexplained input.

\section{Summary}
\label{c06:sec:d96-summary}

This chapter has established the concrete structure of the inevitable spectrum. The D96
spectrum is the Laplacian eigenspectrum of the attractor (Theorem~\ref{c06:thm:d96-derived}):
one zero mode, the background that Difference is counted against
(Theorem~\ref{c06:thm:zero-mode}), and $95$ positive modes (Theorem~\ref{c06:thm:95-modes}).
Its moments $\Sigma m=95$, $\Sigma\sqrt{m}=64.08$, $\Sigma m^2=229$ form the access-count
ladder (Theorem~\ref{c06:cor:moment-ladder}); its constants $\mathrm{occMom}=1900.25$ and
$\mathrm{span}=6.40$ (Theorem~\ref{c06:thm:occmom}, Theorem~\ref{c06:thm:span}) and its
multiplicities $[42\times2,5,6]$ (Theorem~\ref{c06:thm:multiplicity}) are all derived from
the spectrum, which is itself derived from the actualization attractor
(Corollary~\ref{c06:cor:constants-derived}). The spectrum is therefore the substrate of the
organization structure (Theorem~\ref{c06:thm:organization-derived}); the following chapters
read the organization structure and then the physics sectors from this one spectrum.

\begin{thebibliography}{99}
\bibitem{c06:QG155} AT-QG Phase 155, \emph{Symmetry Origin} (Z2 doublet structure), Docs/Research.
\bibitem{c06:QG156} AT-QG Phase 156, \emph{Unified Spectral Access Law}, Docs/Research.
\bibitem{c06:QG157} AT-QG Phase 157, \emph{Effective Access Counts}, Docs/Research.
\bibitem{c06:QG158} AT-QG Phase 158, \emph{Moment Orders}, Docs/Research.
\bibitem{c06:QG159} AT-QG Phase 159, \emph{D96 Selection Origin}, Docs/Research.
\bibitem{c06:QG295} AT-QG Phase 295, \emph{Spectrum Necessity Audit}, Docs/Research.
\bibitem{c06:MONO007} AT-MONO007, \emph{Referee-Readiness Resolution}, Docs/Zenodo/Monograph.
\end{thebibliography}
